\documentclass[11pt,a4paper]{article}
\usepackage[utf8]{inputenc}
\usepackage[T1]{fontenc}
\usepackage[brazil]{babel}
\usepackage[margin=2.2cm]{geometry}
\usepackage{listings}
\usepackage{booktabs}
\usepackage{hyperref}
\setlength{\emergencystretch}{3em}
\sloppy

\lstset{basicstyle=\ttfamily\footnotesize, frame=single, breaklines=true, columns=fullflexible,
  keepspaces=true, literate={á}{{\'a}}1 {ã}{{\~a}}1 {ç}{{\c{c}}}1 {é}{{\'e}}1 {ê}{{\^e}}1}
\hypersetup{colorlinks=true, linkcolor=black, urlcolor=blue}
\setlength{\parskip}{0.3em}

\title{Trabalhos Práticos 1 e 2: Analisador Léxico de Micro C e\\Gerador de Analisadores Sintáticos Descendentes Recursivos a partir de EBNF}
\author{Rafael Torres Nantes \and Victor Hugo de Oliveira Arevalos}
\date{Compiladores I -- FACOM/UFMS \\ 16 de novembro de 2026}

\begin{document}
\maketitle

\section{Introdução e objetivo}

Este relatório descreve os dois trabalhos práticos da disciplina de Compiladores I (UFMS), que
vivem no mesmo repositório porque o segundo usa o primeiro. O \textbf{TP1} é o analisador léxico
da linguagem Micro C, escrito em Flex (pasta \texttt{microc/}). O \textbf{TP2} é uma ferramenta
escrita em C padrão ANSI (C89/C90) que lê uma gramática em EBNF e gera o código-fonte de um
analisador sintático descendente recursivo, com um procedimento \texttt{parse\_A()} para cada não
terminal (pasta \texttt{src/}). O gerador verifica a gramática antes de produzir código, calcula
os conjuntos \emph{anulável}, FIRST e FOLLOW, confirma a condição LL(1) e emite diagnósticos com
arquivo, linha e coluna. No modo Micro C, o analisador gerado pede os tokens ao scanner do TP1,
de modo que a gramática da Micro C do Apêndice B do enunciado é processada de ponta a ponta. A
Seção~\ref{sec:tp1} descreve o TP1 e o que mudou nele para servir ao TP2; as seções seguintes
descrevem o gerador.

A notação de entrada tem duas seções. A seção \texttt{TOKENS} declara terminais por expressão
regular (e \texttt{SKIP} descarta espaços e comentários); a seção \texttt{RULES} traz as
produções, em que o primeiro não terminal é o símbolo inicial. Valem alternativa (\texttt{|}),
sequência, opcional (\texttt{[ ]}), repetição (\texttt{\{ \}}), agrupamento (\texttt{( )}),
literais entre aspas e comentários \texttt{(* ... *)}. Terminais nomeados usam só maiúsculas e
não terminais começam com minúscula. O trabalho foi construído sobre o código-base fornecido
pelo professor, cujos trechos deixados para o aluno foram todos implementados.

\section{TP1: analisador léxico de Micro C}
\label{sec:tp1}

\subsection{Tokens e regras}

O scanner está em \texttt{microc/microc.flex}, e a enumeração \texttt{TokenType} segue a seção
2.3.1 do livro-texto, em \texttt{microc/tokens.h}:

\begin{center}
\begin{tabular}{@{}lp{10.5cm}@{}}
\toprule
Classe & Tokens \\
\midrule
Controle & \texttt{UNDEF} (erro léxico), \texttt{ID}, \texttt{END\_OF\_FILE} \\
Constantes & \texttt{INTEGERCONST}, \texttt{CHARCONST}, \texttt{STRINGCONST} \\
Aritméticos & \texttt{PLUS}, \texttt{MINUS}, \texttt{MUL}, \texttt{DIV}, \texttt{MOD} \\
Relacionais e lógicos & \texttt{EQ}, \texttt{NEQ}, \texttt{LT}, \texttt{GT}, \texttt{LEQ},
  \texttt{GEQ}, \texttt{AND}, \texttt{OR}, \texttt{NOT} \\
Pontuação & \texttt{ASSIGN}, \texttt{SEMICOLON}, \texttt{COMMA}, parênteses, chaves e colchetes \\
Reservadas & \texttt{MAIN}, \texttt{IF}, \texttt{ELSE}, \texttt{FOR}, \texttt{RETURN},
  \texttt{INT}, \texttt{CHAR}, \texttt{PRINT} \\
\bottomrule
\end{tabular}
\end{center}

As regras aproveitam duas propriedades do Flex: vence o \textbf{maior casamento} e, em empate, a
regra escrita primeiro. Por isso os operadores de dois caracteres (\texttt{==}, \texttt{!=},
\texttt{<=}, \texttt{>=}, \texttt{\&\&}, \texttt{||}) vêm antes dos de um caractere com o mesmo
prefixo. Identificadores casam \texttt{\{LETRA\}\{ALFANUM\}*} e são comparados com as palavras
reservadas depois do casamento, o que evita uma regra por palavra. Inteiros negativos
(\texttt{-}\{DIGIT\}+) são uma constante só. Espaços e comentários de linha (\texttt{//}) são
descartados, e cada \texttt{\textbackslash n} incrementa o contador de linha. O comentário de
bloco usa a condição de início \texttt{COMMENT}, que consome tudo até \texttt{*/} e continua
contando linhas. Constantes de caractere e de string aceitam sequências de escape
(\texttt{\textbackslash.}). Os lexemas de identificadores e constantes são guardados numa tabela
de strings, e o token carrega tipo, lexema e linha.

\subsection{Erros léxicos}

Toda entrada casa alguma regra: a última, \texttt{.}, transforma qualquer caractere não previsto
em \texttt{UNDEF}. Cada erro devolve \texttt{UNDEF} com a mensagem em \texttt{error\_msg}, e a
análise continua no caractere seguinte.

\begin{center}
\begin{tabular}{@{}ll@{}}
\toprule
Situação & Mensagem \\
\midrule
Fim do arquivo dentro de \texttt{/* ... */} & \texttt{EOF em comentario} \\
\texttt{*/} fora de um comentário & \texttt{Comentario nao iniciado} \\
Caractere vazio, sem fechar ou com mais de um símbolo & \texttt{Constante de caractere invalida} \\
String com caractere nulo & \texttt{String contem caractere nulo} \\
Quebra de linha antes de fechar a string & \texttt{String nao terminada} \\
Fim do arquivo antes de fechar a string & \texttt{EOF em string} \\
Caractere fora do alfabeto & o próprio caractere \\
\bottomrule
\end{tabular}
\end{center}

\subsection{Mudanças no TP1 para servir ao TP2}

A entrega original do TP1 era um programa autônomo: o \texttt{microc.flex} com \texttt{main()},
que imprimia a lista de tokens de um arquivo. No TP2, o scanner precisa ser chamado pelo
analisador gerado, um token por vez, sem conflitar com o resto do programa. As regras léxicas
continuaram as mesmas; mudou a forma de ligar o scanner ao resto:
\begin{itemize}
\item \texttt{\%option prefix="microc\_yy" noyywrap noinput nounput}: o scanner passa a se chamar
  \texttt{microc\_yylex()} e não colide com o léxico que o gerador escreve para as outras
  gramáticas. \texttt{noinput} e \texttt{nounput} retiram funções não usadas, que gerariam aviso
  com \texttt{-Wall}.
\item \texttt{enum TokenType} e \texttt{YYSTYPE} saíram do \texttt{.flex} para
  \texttt{microc/tokens.h}, que o gerador lê para descobrir os nomes dos terminais.
\item O \texttt{main()} de teste só é compilado com \texttt{-DMICROC\_TESTE\_LEXICO}; sem a
  macro, o único \texttt{main} é o do analisador gerado.
\item A \texttt{yywrap()} manual deu lugar a \texttt{noyywrap}, e a nova
  \texttt{microc\_libera\_tabela()} libera a tabela de strings no fim da análise.
\item O adaptador novo \texttt{microc/scanner.c} expõe três funções: \texttt{scanner\_iniciar},
  \texttt{proximo\_token} e \texttt{scanner\_finalizar}. A segunda monta um \texttt{Token} com o
  lexema vindo da tabela de strings e, para \texttt{UNDEF}, com a mensagem de erro do TP1.
\end{itemize}
A entrega autônoma (\texttt{tests/test.mc} e \texttt{leiame.txt}) saiu da árvore do repositório e
continua no histórico do git.

\section{TP2: arquitetura do gerador}

O gerador é um \emph{pipeline} de cinco etapas, cada uma em seu módulo de \texttt{src/}:

\begin{center}
\small
\begin{tabular}{@{}rll@{}}
\toprule
Etapa & Módulo & Produto \\
\midrule
1. Leitura da EBNF & \texttt{ebnf\_lexer.c}, \texttt{ebnf\_parser.c} & árvore \texttt{Expr} por regra \\
2. Verificação semântica & \texttt{verificacao.c} & erros e avisos de definição e alcance \\
3. Conjuntos & \texttt{conjuntos.c}, \texttt{bitset.c} & anulável, FIRST, FOLLOW \\
4. Verificação LL(1) & \texttt{ll1.c} & erros e avisos de conflito \\
5. Geração de código & \texttt{codegen.c}, \texttt{automato.c} & \texttt{prefixo.c} e \texttt{prefixo.h} \\
\bottomrule
\end{tabular}
\end{center}

O \texttt{main.c} chama as etapas nessa ordem e para na primeira que reporta erro, antes de abrir
qualquer arquivo de saída. Assim, uma gramática recusada nunca deixa um \texttt{.c} pela metade.

\subsection{Estruturas de dados}

Cada regra é guardada como uma árvore \texttt{Expr} (\texttt{src/gramatica.h}), com sete tipos de
nó:

\begin{center}
\begin{tabular}{@{}lll@{}}
\toprule
EBNF & Nó & Filhos \\
\midrule
terminal \texttt{t} & \texttt{E\_TERM} & índice em \texttt{g->terms} \\
não terminal \texttt{B} & \texttt{E\_NAOTERM} & índice em \texttt{g->nts} \\
\texttt{a b c} & \texttt{E\_SEQ} & \texttt{a}, \texttt{b}, \texttt{c} \\
\texttt{a | b | c} & \texttt{E\_ALT} & \texttt{a}, \texttt{b}, \texttt{c} \\
\texttt{[ a ]} & \texttt{E\_OPC} & \texttt{a} \\
\texttt{\{ a \}} & \texttt{E\_REP} & \texttt{a} \\
sequência vazia & \texttt{E\_VAZIO} & -- \\
\bottomrule
\end{tabular}
\end{center}

Cada nó guarda também a linha e a coluna de origem, usadas nas mensagens, e os seus próprios
conjuntos anulável, FIRST e FOLLOW, preenchidos na etapa 3. Os terminais, inclusive os literais
entre aspas, ganham um índice numa tabela única, e os conjuntos são vetores de bits
(\texttt{bitset.c}) indexados por ele. União, interseção e teste de vazio custam uma passada pelo
vetor, o que mantém barato o ponto fixo, que repete essas operações até estabilizar.

\subsection{Léxico e árvore da EBNF}

O analisador da própria notação é descendente recursivo: \texttt{parse\_expr} trata as
alternativas, \texttt{parse\_seq} as sequências e \texttt{parse\_fator} os fatores
(terminal, não terminal, literal, \texttt{[ ]}, \texttt{\{ \}} e \texttt{( )}). Sequência e
alternativa com um só filho são devolvidas diretamente, a sequência vazia vira um nó
\texttt{E\_VAZIO} e o agrupamento não cria nó. A árvore resultante, um \texttt{Expr} por regra,
é a estrutura consumida por todas as etapas seguintes.

\subsection{Análise semântica}

A etapa 2 detecta os erros de definição e de alcance exigidos pelo enunciado: não terminal usado
e não definido, não terminal definido mais de uma vez e terminal nomeado não declarado em
\texttt{TOKENS}. Em seguida vêm três análises sobre o grafo da gramática:
\begin{itemize}
\item \textbf{Inalcançáveis} (aviso): busca em profundidade a partir do símbolo inicial.
\item \textbf{Improdutivos} (erro): ponto fixo que marca como produtivo o não terminal que
  deriva alguma cadeia de terminais; os não marcados ao final são improdutivos.
\item \textbf{Recursão à esquerda} (erro): fecho transitivo (Warshall) do grafo \emph{pode
  começar com}, que atravessa prefixos anuláveis. Assim a recursão indireta e a recursão por
  prefixo anulável aparecem como ciclo; refazer apenas a linha do próprio não terminal separa o
  caso direto do indireto, para a mensagem ser precisa. Por depender do anulável, essa função
  de \texttt{verificacao.c} é chamada no início da etapa 3, antes do FIRST, cujo ponto fixo
  nunca terminaria numa regra recursiva à esquerda.
\end{itemize}

\subsection{Anulável, FIRST e FOLLOW}

Os três conjuntos são calculados por ponto fixo, com uma passada por regra até que nada mude,
para cada não terminal e para cada subexpressão da árvore. Os conjuntos são bitsets
(\texttt{bitset.c}). Anulável e FIRST seguem o tipo do nó:

\begin{center}
\begin{tabular}{@{}lll@{}}
\toprule
Nó & Anulável & FIRST \\
\midrule
\texttt{E\_TERM} \texttt{t} & não & $\{t\}$ \\
\texttt{E\_NAOTERM} \texttt{B} & anulável(B) & FIRST(B) \\
\texttt{E\_SEQ} $a_1 \dots a_n$ & todos os $a_i$ anuláveis & $\bigcup$ FIRST($a_i$) enquanto o prefixo anterior for anulável \\
\texttt{E\_ALT} & algum filho anulável & $\bigcup$ FIRST dos filhos \\
\texttt{E\_OPC}, \texttt{E\_REP} & sim & FIRST do filho \\
\texttt{E\_VAZIO} & sim & $\emptyset$ \\
\bottomrule
\end{tabular}
\end{center}

\noindent O valor de um não terminal é o da raiz da sua regra; como uma regra pode citar outra,
ainda não calculada, a passada se repete até que nenhum bit mude. O FOLLOW percorre cada sequência de trás para a frente, acumulando o que
pode vir depois de cada filho; no corpo de uma repetição \texttt{\{ a \}}, \texttt{follow(a)}
recebe também \texttt{FIRST(a)}, pois depois de uma iteração pode vir outra. O FOLLOW do símbolo
inicial contém o marcador de fim de entrada \texttt{\$}. A opção \texttt{--sets} imprime a
tabela.

\subsection{Verificação LL(1)}

O módulo \texttt{ll1.c} aplica as condições do enunciado a cada alternativa, opcional e
repetição. São \textbf{erros}, e nenhum código é gerado: FIRST não disjuntos entre alternativas;
mais de uma alternativa anulável; alternativa anulável cujo FOLLOW intersecta o FIRST de outra;
e corpo de \texttt{[ ]} ou \texttt{\{ \}} anulável. É \textbf{aviso} a interseção FIRST $\cap$
FOLLOW em \texttt{[ ]} ou \texttt{\{ \}}, resolvida de forma gulosa (o caso do \emph{else}
pendente).

\subsection{Geração de código}

Cada regra vira uma função C, com a regra EBNF de origem em comentário. A sequência é emitida em
ordem; a alternativa vira \texttt{if / else if}, com a alternativa anulável no \texttt{else}
final no lugar do erro; \texttt{[ ]} vira \texttt{if} e \texttt{\{ \}} vira \texttt{while}. Os
conjuntos de decisão saem como \texttt{static const int conj\_N[] = \{ ..., -1 \};} e o erro
sintático lista os tokens esperados, obtidos de FIRST. A regra \texttt{term} da gramática de
expressões, por exemplo, vira:

\begin{lstlisting}[language=C]
/* term = factor { ( '*' | '/' ) factor } ; */
void parse_term(void)
{
    arvore_entra("term");
    parse_factor();
    while (tok.tipo == T_VEZES || tok.tipo == T_BARRA) {
        if (tok.tipo == T_VEZES) {
            match(T_VEZES);
        } else if (tok.tipo == T_BARRA) {
            match(T_BARRA);
        } else {
            erro_sintatico(conj_1);
        }
        parse_factor();
    }
    arvore_sai();
}
\end{lstlisting}

\noindent A condição do \texttt{while} é o FIRST do corpo da repetição; \texttt{match} consome o
token esperado ou para com erro; \texttt{conj\_1} é \texttt{\{ T\_VEZES, T\_BARRA, -1 \}}, a
lista que aparece na mensagem. Os nomes \texttt{T\_...} vêm do \texttt{enum} gravado em
\texttt{prefixo.h}, e literais recebem nomes legíveis (\texttt{'*'} vira \texttt{T\_VEZES},
\texttt{'('} vira \texttt{T\_ABRE\_PAR}).

\paragraph{Opções do analisador gerado.} \texttt{arvore\_entra} e \texttt{arvore\_sai} só imprimem
quando o analisador foi gerado com \texttt{--tree} e executado com \texttt{--tree}: cada chamada de
procedimento é um nível de recuo, e cada token consumido é uma folha. Com \texttt{--tokens
arquivo}, os tokens vêm de um arquivo com uma linha por token, no formato do livro-texto
(\texttt{Token: tipo = 34, lexema = 'int', linha = 1}), em vez do léxico. Esse modo isola o
analisador sintático: se a leitura direta e a leitura do arquivo derem saídas diferentes, o
defeito está no léxico ou no adaptador, e não nas regras.

\paragraph{Memória.} O gerador libera a gramática, a tabela de tokens da Micro C e o texto lido
antes de sair. No analisador gerado, a função \texttt{limpar()} fecha a entrada e libera o buffer
e o scanner; ela é registrada com \texttt{atexit} e roda tanto na aceitação quanto em qualquer
saída por erro léxico ou sintático. Essa correção veio da revisão final: antes dela, toda entrada
rejeitada terminava com arquivo aberto e memória alocada.

\paragraph{Léxico próprio.} Para as gramáticas com seção \texttt{TOKENS}, o gerador emite também
o analisador léxico. Cada expressão regular é convertida em AFN (construção de Thompson), em AFD
(construção de subconjuntos) e em AFD mínimo (algoritmo de Moore, \texttt{automato.c}). O AFD
mínimo é emitido como um \texttt{switch} por estado, com transições agrupadas em faixas de
caracteres. Vale o maior prefixo e, em empate, o literal vence o terminal nomeado. A regex que
aceita a cadeia vazia é recusada.

\paragraph{Integração com o TP1.} Com \texttt{--scanner microc} (declaração
\texttt{TOKENS EXTERN "tokens.h";}), o analisador chama o scanner flex do TP1 sob demanda. Os
nomes dos tokens vêm do \texttt{enum TokenType} lido de \texttt{tokens.h} pelo módulo
\texttt{microc\_mapa.c}. O adaptador em \texttt{microc/scanner.c} obtém o lexema da tabela de
strings do TP1 e a mensagem do token \texttt{UNDEF} de \texttt{error\_msg}; o fim do arquivo
vira \texttt{END\_OF\_FILE}. Ao final, \texttt{scanner\_finalizar()} libera a tabela.

\section{Decisões de projeto e tratamento de erros}

\paragraph{Recursão à esquerda.} Em vez de procurar só o ciclo direto \texttt{A $\to$ A ...},
o fecho transitivo do grafo \emph{pode começar com} cobre a recursão indireta e a que passa por
prefixo anulável, como em \texttt{a = [","] b ; b = a ;}. Um analisador descendente recursivo
geraria laço infinito nesses casos, por isso são erros fatais (gramática
\texttt{erro\_recursao\_indireta\_anulavel}).

\paragraph{Improdutivos.} Um não terminal como \texttt{laco = "x" laco ;}, sem caso base, nunca
deriva uma cadeia de terminais. É erro, e não aviso, porque o procedimento gerado jamais
terminaria com sucesso.

\paragraph{Repetição e opcional anuláveis.} O corpo \texttt{\{ [ "x" ] \}} faria o \texttt{while}
gerado girar sem consumir token. A verificação LL(1) rejeita corpo anulável em \texttt{[ ]} e
\texttt{\{ \}}.

\paragraph{\emph{Else} pendente.} A produção do comando \texttt{if} da Micro C tem a forma
\texttt{"if" "(" expr ")" comando [ "else" comando ]}. O terminal \texttt{else} está em FIRST do
opcional e também em FOLLOW de \texttt{comando\_if}, logo a gramática não é estritamente LL(1).
O gerador não rejeita: emite o aviso exigido e resolve de forma gulosa, ligando o \texttt{else} ao
\texttt{if} mais próximo, que é a semântica desejada. O aviso que aparece na compilação é

\begin{lstlisting}
microc.ebnf:19:57: aviso: conflito LL(1) em comando_if ... {'else'}
\end{lstlisting}

\paragraph{Fatoração à esquerda.} A Micro C do Apêndice B exige fatoração em \texttt{fator} e
\texttt{atribuicao\_interna}, para que as alternativas tenham FIRST disjuntos. A fatoração foi
feita na própria gramática de entrada, mantendo o gerador simples.

\paragraph{Código gerado seguro.} Literais que contêm \texttt{*/} ou \texttt{/*} quebrariam o
comentário C que documenta a regra. O código fornecido foi corrigido para emitir esses pares com
um espaço no meio, o que evita erro de compilação e o aviso \texttt{-Wcomment}.

\paragraph{Mensagens.} Todo diagnóstico segue o formato \texttt{arquivo:linha:coluna: erro:} ou
\texttt{aviso:}. Erro encerra com código 1, antes de qualquer geração; aviso não impede a
geração. Arquivo de gramática inexistente também encerra o gerador com código 1; o analisador
gerado devolve 2 quando não consegue abrir o programa de entrada. Expressões regulares de
\texttt{TOKENS} e \texttt{SKIP} são compiladas já na etapa 2, de modo que uma regex inválida
(intervalo invertido, ou que aceita a cadeia vazia) é erro com posição e não deixa arquivo
gerado, mesmo quando o gerador roda só com \texttt{--sets}.

\section{Testes e resultados}

Os resultados abaixo foram medidos em 16/11/2026 sobre uma cópia limpa do repositório, em dois
ambientes: Ubuntu 24.04 (gcc 13.3.0, flex 2.6.4, valgrind 3.22.0), o sistema em que o trabalho é
corrigido, e Windows 11 com o MSYS2 (gcc 15.3.0).

\begin{center}
\small
\begin{tabular}{@{}lp{7.6cm}r@{}}
\toprule
Verificação & Resultado & \\
\midrule
\texttt{make clean \&\& make} & sem avisos do compilador; só o diagnóstico exigido do \emph{else} pendente & saída 0 \\
\texttt{TODO(aluno)} & \texttt{grep -rn "TODO(aluno)" src microc} & nenhum \\
\texttt{make test} & testes do professor e do grupo, com \texttt{make microc} & 35 de 35 \\
Código gerado & \texttt{expr}, \texttt{json}, \texttt{comandos} e Micro C com \texttt{gcc -ansi -pedantic -Wall -Wextra} & sem avisos \\
\texttt{valgrind} & \texttt{--leak-check=full}: o gerador nas 10 gramáticas e cada analisador em todos os seus programas, com e sem \texttt{--tree} & 60 execuções, 0 erros \\
\texttt{--tokens} & nos 6 programas de Micro C, a saída e o código de saída com \texttt{--tokens} são iguais aos da leitura direta & 6 de 6 \\
\bottomrule
\end{tabular}
\end{center}

O \texttt{make test} executa \texttt{testes/executar.sh}, que roda o gerador com \texttt{--sets}
em cada gramática e cada analisador em cada programa, e compara a saída, inclusive a linha final
\texttt{status 0} ou \texttt{status 1}, com o arquivo esperado por \texttt{diff}. O arquivo de
tokens do teste de \texttt{--tokens} foi gravado por um programa auxiliar que chama
\texttt{proximo\_token()} até \texttt{END\_OF\_FILE}.

\subsection{Gramáticas}

Foram escritas dez gramáticas de teste em \texttt{testes/gramaticas/}: quatro válidas e seis com
erro (o enunciado exige ao menos três válidas, uma delas a Micro C, e duas com erro).

\begin{center}
\small
\begin{tabular}{@{}llp{8.2cm}@{}}
\toprule
Gramática & Tipo & O que exercita \\
\midrule
\texttt{expr} & válida & expressões aritméticas, precedência por níveis \\
\texttt{microc} & válida & Micro C com o scanner do TP1; \emph{else} pendente \\
\texttt{comandos} & válida & \texttt{if}/\texttt{while}, blocos, \emph{else} pendente \\
\texttt{json} & válida & não terminais anuláveis, listas e objetos \\
\texttt{conflito\_ll1} & erro & FIRST não disjuntos em alternativas \\
\texttt{recursao\_esquerda} & erro & recursão à esquerda direta \\
\texttt{erro\_semantica} & erro & não terminal indefinido, terminal não declarado, inalcançável \\
\texttt{erro\_improdutivo} & erro & \texttt{laco} nunca termina \\
\texttt{erro\_recursao\_indireta\_anulavel} & erro & \texttt{a = [","] b ; b = a ;} \\
\texttt{erro\_repeticao\_anulavel} & erro & \texttt{\{ [ "x" ] \}} \\
\bottomrule
\end{tabular}
\end{center}

\subsection{Programas aceitos e rejeitados}

Cada gramática válida tem ao menos três programas aceitos e três rejeitados, o mínimo exigido:
\texttt{comandos}, \texttt{json} e \texttt{microc} têm 3 e 3, e \texttt{expr} tem 3 e 4, o que
dá 25 programas em \texttt{testes/programas/}. As saídas esperadas do \texttt{--sets} estão em
\texttt{testes/esperado/}; as dos programas, ao lado de cada um, em arquivos \texttt{.esperado}.

\subsection{Robustez}

Além dos 35 testes, o gerador e os analisadores gerados foram exercitados com entradas no limite:
gramáticas vazias, literais ou expressões regulares sem fechar, nomes reservados do C, arquivos
com CRLF, parênteses aninhados em excesso e bytes aleatórios. O critério é que o gerador e os
analisadores terminem sempre com 0 ou 1 (2 para arquivo inexistente), sem travar nem cair, e que
todo código gerado compile com \texttt{-Werror}. Esse exercício encontrou o defeito do literal
com \texttt{*/} descrito acima.

\subsection{Exemplos de saída}

Para a entrada \texttt{2 * x}, o analisador de expressões gerado com \texttt{--tree} imprime a
árvore de derivação e aceita; para \texttt{2 * (x + )}, aponta a posição e os tokens esperados:

\begin{lstlisting}
$ ./build/expr_bin --tree exemplo.txt
expr
  term
    factor
      NUM '2'
    '*'
    factor
      ID 'x'
entrada aceita
$ ./build/expr_bin erro.txt
linha 1, coluna 10: encontrado ')', esperado um de {NUM, ID, '('}
\end{lstlisting}

\noindent No modo Micro C, o scanner do TP1 só registra a linha, então as mensagens não têm
coluna, e o erro léxico traz a mensagem do TP1:

\begin{lstlisting}
linha 3: encontrado ';', esperado um de {'(', ID, INTEGERCONST, CHARCONST, '!'}
linha 2: erro lexico: String nao terminada
\end{lstlisting}

\subsection{Revisão final do código}

Depois de os testes passarem, o código foi revisado função por função, e duas mudanças entraram:
a liberação de memória do analisador gerado por \texttt{atexit}, descrita na Seção 3, e a
remoção de trechos redundantes em \texttt{automato.c}, \texttt{conjuntos.c} e
\texttt{verificacao.c} (41 linhas removidas e 16 acrescentadas, sem mudar nenhuma saída). Os 35
testes foram repetidos depois de cada mudança.

\section{Limitações}

\begin{itemize}
\item \textbf{Conflitos não são corrigidos automaticamente.} O gerador diagnostica, mas não
  fatora nem reescreve a gramática; a fatoração da Micro C é feita à mão.
\item \textbf{Sem recuperação de erro.} O analisador gerado para no primeiro erro sintático e
  lista os tokens esperados; não continua a análise depois dele.
\item \textbf{LL(1) com resolução gulosa.} Só o conflito de \texttt{[ ]} e \texttt{\{ \}} é
  tolerado, com aviso; qualquer outro conflito impede a geração.
\item \textbf{Sem coluna no modo Micro C.} O scanner do TP1 só conta linhas, e as mensagens desse
  modo trazem apenas a linha.
\item \textbf{Expressões regulares restritas.} O léxico próprio aceita caractere, escape
  \texttt{\textbackslash x}, \texttt{.}, classes \texttt{[abc]}, \texttt{[a-z]} e \texttt{[\^{}...]},
  agrupamento, \texttt{|}, \texttt{*}, \texttt{+} e \texttt{?}; não há quantificadores
  \texttt{\{n,m\}} nem âncoras.
\item \textbf{Ambiente.} Testado no gcc 13.3.0 (Ubuntu 24.04) e no gcc 15.3.0 (MSYS2); o valgrind
  só roda no Linux.
\end{itemize}

\section{Uso de IA}

O código do gerador foi escrito pelos autores. Uma ferramenta de IA (Claude, da Anthropic) foi
usada como apoio na fase de testes -- sugestão de casos e conferência dos resultados -- e na
redação e revisão deste relatório.

\section{Como compilar e executar}

São necessários \texttt{gcc}, \texttt{make}, \texttt{flex} e \texttt{diff}. No Windows, o
\emph{toolchain} do MSYS2 precisa estar no PATH (ou a sessão deve ser aberta pelo \emph{shell}
do MSYS2); sem isso o \texttt{make} não é encontrado. Na pasta do projeto:

\begin{lstlisting}
make                          # compila ./gerador e build/microc/parser_microc
make test                     # testes de testes/ (35)
make clean
\end{lstlisting}

Para gerar e usar um analisador:

\begin{lstlisting}
./gerador testes/gramaticas/expr.ebnf -o build/expr --tree
gcc -ansi -pedantic -Wall -Wextra -o build/expr_bin build/expr.c
./build/expr_bin --tree programa.txt
./build/expr_bin --tokens tokens.txt
\end{lstlisting}

As opções do gerador são: \texttt{-o prefixo} (grava \texttt{prefixo.c} e \texttt{prefixo.h}),
\texttt{--sets} (imprime anulável, FIRST e FOLLOW), \texttt{--debug} (imprime a gramática lida),
\texttt{--tree} (inclui a impressão da árvore de derivação) e \texttt{--scanner microc} (usa o
scanner do TP1). O código, o scanner adaptado em \texttt{microc/}, os testes e o registro de uso
de IA (\texttt{USO\_IA.md}) estão no repositório do trabalho.

\end{document}
